\section{Background and Related Work}
\label{sec:related}

Corporate lobbying has demonstrable consequences for climate policy outcomes.
Brulle~\cite{brulle2018climate} maps US climate-lobbying expenditure across industry sectors, and Meng and Rode~\cite{meng2019social} estimate that lobbying lowered the probability of enacting the Waxman--Markey Cap-and-Trade Bill (a crucial climate action bill in the US) by 13 percentage points, with an estimated welfare loss of roughly \$60 billion.
More recently, Leippold, Sautner and Yu~\cite{leippold2024corporate} document a \emph{camouflaging trend}, whereby firms increasingly conceal climate lobbying behind opaque legislative references, while Baehr, Bare and Heddesheimer~\cite{baehr2026climate} show across more than 11{,}000 firms that earnings-call \emph{climate exposure} systematically predicts whether, how much, and where a firm lobbies.
Understanding lobbying therefore requires more than aggregate spending figures: it requires modelling the underlying network of actors, funders, and affiliations through which climate-policy influence is organised and exerted.

\begin{sloppypar}Investigative platforms partly address the documentation challenge.
DeSmog~\cite{desmog} profiles individuals and organisations involved in denial and lobbying; LobbyMap~\cite{lobbymap} tracks corporate engagement with climate policy; Skeptical Science~\cite{skepticalscience} curates statements attributed to prominent misinformation actors; and Oreskes and Conway~\cite{oreskes2011merchants} analyse strategies used to manufacture doubt about scientific consensus.
These resources are invaluable but heterogeneous, siloed, and not interlinked in any form amenable to computational analysis.
Concurrent efforts such as the UCL \emph{AI-Powered Climate Lobbying Database}~\cite{uclclimatelobbying} reflect the recognised need for better support, automating text extraction and network analysis using large language models.\end{sloppypar}

In the Semantic Web community, CimpleKG~\cite{burel2024cimplekg} provides a continuously updated knowledge graph on misinformation and fact-checks across domains, currently being extended within the ClimateSense project~\cite{climatesense} with climate-specific classifiers, while ClimaFactsKG~\cite{burel2025climafactskg} structures scientific evidence countering climate misinformation; both, however, target claim verification rather than lobbying actor networks.
Knowledge graphs have also been adopted as infrastructure for investigative journalism~\cite{opdahl2022semantic,gallofre2022jkp,berven2020kg,alexander2023manifesto,motta2025epistemology}, but with a focus on news content and event extraction rather than political and corporate actors.
Adjacent Linked Open Data efforts cover political discourse, including the European Parliament debates~\cite{van2016debates} and the digitised Canadian parliamentary record~\cite{beelen2017digitization}, but model what political actors say in formal proceedings rather than the off-stage networks of funding, affiliation, and personnel through which lobbying influence operates.

Constructing such a graph requires entity resolution and entity linking, both well-studied: recent surveys cover end-to-end ER pipelines~\cite{christophides2020er} and neural EL methods~\cite{sevgili2022neural}, with state-of-the-art systems combining candidate retrieval, learned similarity, and graph-aware disambiguation~\cite{ayoola2022refined}.
Large language models have more recently been applied directly to entity matching, often matching or exceeding supervised baselines without task-specific training~\cite{peeters2025llmem}.
LLM-driven approaches fit the climate-lobbying setting: source documents demand contextual disambiguation, and many mentions are long-tail entities absent from general-purpose knowledge bases.
LACK's hybrid pipeline (Section~\ref{sec:extraction}) reflects these constraints, combining knowledge-base candidate retrieval, LLM-based disambiguation, and a web search fallback.

What is missing is a machine-readable representation of the actors and relationships at the core of climate lobbying.
LACK fills this gap as a Wikidata-linked knowledge graph built from investigative sources through an evaluated extraction pipeline and published as Linked Open Data.